% 模板 10：管道分段 + 并行分支 + 汇聚（数据科学/竞赛风格）
% 统一精致配色：蓝色主节点+teal子节点+微阴影
% 编译命令：xelatex demo_roadmap_research_pipeline.tex
\documentclass[border=5pt]{standalone}
\usepackage{ctex}
\usepackage{tikz}
\usetikzlibrary{arrows.meta, positioning, calc, shadows}

\begin{document}
\begin{tikzpicture}[
    dashbox/.style={rectangle, rounded corners=4pt,
        draw=gray!40, dashed, line width=0.7pt, fill=gray!2},
    main/.style={rectangle, rounded corners=3pt,
        minimum width=3.8cm, minimum height=1.4cm, align=center,
        draw=blue!80, line width=0.7pt, fill=blue!6, font=\small,
        drop shadow={opacity=0.15, shadow xshift=0.5pt, shadow yshift=-0.5pt}},
    sub/.style={rectangle, rounded corners=2pt,
        minimum width=2.2cm, minimum height=0.5cm, align=center,
        draw=teal!70, line width=0.5pt, fill=white, font=\scriptsize},
    bigarrow/.style={-stealth, line width=1.4pt, color=blue!70},
    smarrow/.style={-stealth, line width=0.5pt, color=gray!50},
    label/.style={font=\small\bfseries, color=black!70},
    branch/.style={-stealth, line width=0.6pt, color=blue!50, densely dashed},
    merge/.style={circle, fill=blue!65, minimum size=11pt,
        font=\tiny\bfseries, text=white, inner sep=0pt},
]

% ====== Stage 1: 问题定义 ======
\node[dashbox, minimum width=14cm, minimum height=2.8cm] (box1) at (0, 0) {};
\node[label, anchor=north west] at (box1.north west) {\scriptsize 阶段一：问题定义与数据获取};
\node[main] (p1a) at (-4, 0) {\textbf{问题建模}\\[2pt]{\scriptsize 目标函数/约束}};
\node[main] (p1b) at (0, 0) {\textbf{数据采集}\\[2pt]{\scriptsize 多源异构数据}};
\node[main] (p1c) at (4, 0) {\textbf{质量评估}\\[2pt]{\scriptsize 完整性/一致性}};
\draw[smarrow] (p1a) -- (p1b); \draw[smarrow] (p1b) -- (p1c);

\draw[bigarrow] (0, -1.7) -- (0, -2.5);

% ====== Stage 2: 特征工程 ======
\node[dashbox, minimum width=14cm, minimum height=2.8cm] (box2) at (0, -4.2) {};
\node[label, anchor=north west] at (box2.north west) {\scriptsize 阶段二：特征工程与预处理};
\node[main] (p2a) at (-4, -4.2) {\textbf{缺失值处理}\\[2pt]{\scriptsize 多重插补/KNN}};
\node[main] (p2b) at (0, -4.2) {\textbf{特征构造}\\[2pt]{\scriptsize 交互项/滞后}};
\node[main] (p2c) at (4, -4.2) {\textbf{降维与选择}\\[2pt]{\scriptsize PCA/LASSO}};
\draw[smarrow] (p2a) -- (p2b); \draw[smarrow] (p2b) -- (p2c);

\draw[bigarrow] (0, -5.9) -- (0, -6.7);

% ====== Stage 3: 模型构建（并行分支） ======
\node[dashbox, minimum width=14cm, minimum height=5.5cm] (box3) at (0, -9.8) {};
\node[label, anchor=north west] at (box3.north west) {\scriptsize 阶段三：模型构建与训练};

\node[font=\scriptsize, color=blue!50] at (0, -7.2) {$\longleftarrow$ 并行建模对比 $\longrightarrow$};

\node[main, minimum width=3.5cm] (p3a) at (-4.2, -8.8) {\textbf{统计/计量模型}\\[2pt]{\scriptsize FE/DID/IV-2SLS}};
\node[main, minimum width=3.5cm] (p3b) at (0, -8.8) {\textbf{机器学习模型}\\[2pt]{\scriptsize XGBoost/RF/SVM}};
\node[main, minimum width=3.5cm] (p3c) at (4.2, -8.8) {\textbf{深度学习模型}\\[2pt]{\scriptsize LSTM/Transformer}};

\draw[branch] (-2, -7.0) -- (p3a.north);
\draw[branch] (0, -7.0) -- (p3b.north);
\draw[branch] (2, -7.0) -- (p3c.north);

\node[merge] (m3) at (0, -10.2) {$\Sigma$};
\draw[smarrow] (p3a.south) |- (m3);
\draw[smarrow] (p3b.south) -- (m3);
\draw[smarrow] (p3c.south) |- (m3);

\node[main, minimum width=5.5cm] (tune) at (0, -11.4) {\textbf{超参数优化与模型选择}\\[2pt]{\scriptsize 网格搜索/贝叶斯优化/交叉验证}};
\draw[smarrow] (m3) -- (tune);

\draw[bigarrow] (0, -12.4) -- (0, -13.2);

% ====== Stage 4: 评估验证 ======
\node[dashbox, minimum width=14cm, minimum height=2.8cm] (box4) at (0, -14.9) {};
\node[label, anchor=north west] at (box4.north west) {\scriptsize 阶段四：评估验证与结果解释};
\node[main] (p4a) at (-4, -14.9) {\textbf{性能评估}\\[2pt]{\scriptsize RMSE/$R^2$/AUC}};
\node[main] (p4b) at (0, -14.9) {\textbf{可解释性分析}\\[2pt]{\scriptsize SHAP/偏依赖图}};
\node[main] (p4c) at (4, -14.9) {\textbf{稳健性检验}\\[2pt]{\scriptsize 消融/敏感性}};
\draw[smarrow] (p4a) -- (p4b); \draw[smarrow] (p4b) -- (p4c);

\draw[bigarrow] (0, -16.6) -- (0, -17.4);

% ====== Stage 5: 结论 ======
\node[dashbox, minimum width=14cm, minimum height=2cm] (box5) at (0, -18.6) {};
\node[label, anchor=north west] at (box5.north west) {\scriptsize 阶段五：结论与决策建议};
\node[main] (p5a) at (-4, -18.6) {\textbf{核心发现}\\[2pt]{\scriptsize 关键变量识别}};
\node[main] (p5b) at (0, -18.6) {\textbf{方案优化}\\[2pt]{\scriptsize 最优策略输出}};
\node[main] (p5c) at (4, -18.6) {\textbf{政策建议}\\[2pt]{\scriptsize 实施路径}};
\draw[smarrow] (p5a) -- (p5b); \draw[smarrow] (p5b) -- (p5c);

\end{tikzpicture}
\end{document}
